% Section IV — Evaluation Methodology
% Draft 2026-10-20. Numbers: eval/ANSWER_EVAL_REPORT.md, eval/RETRIEVAL_REPORT.md.
% References: references.bib (verified 2026-10-20).

\section{Evaluation Methodology}
\label{sec:method}

\subsection{Questions}

We use 50 natural-language questions over ORKG in six types:
topic search (9), method usage (9), dataset search (8), method comparison (8),
claim verification (8) and paper lookup (8); 14 are easy, 20 medium and 16 hard.
Questions with odd IDs form the development split and those with even IDs
the test split (25 each). All retrieval development (strategies, ranking,
query-classification rules) was done on the development split only; we report
results on all 50 questions and on the test split.

\subsection{Relevance Gold Set}

For every question we wrote a grading criterion and graded candidate papers
on a three-level scale: 2 (the paper answers the question, e.g.\ it reports
both compared methods on the named task), 1 (partially relevant, e.g.\ one of
the two methods, or the right method on another task) and 0 (not relevant).
Candidates were pooled from every retrieval variant run during development,
from both systems compared in this paper (including a second embedding model),
and from all papers cited in generated answers. The final set contains 2,633
judgments (757 of grade 2, 890 of grade 1, 986 of grade 0); every paper in the
top 8 of every reported system is judged.

Grades were proposed by an LLM (Claude) from each paper's ORKG record and
reviewed by the author. To estimate their reliability, the author re-graded a
stratified sample of 120 question--paper pairs (40 per grade) without seeing
the proposed grades. Agreement was 0.85 (Cohen's $\kappa = 0.78$, quadratic-weighted
$\kappa = 0.85$; $\kappa = 0.91$ for the binary decision ``grade 2 or not''),
and $\kappa = 0.91$ on the pairs retrieved only by the dense baseline, so the
added grades do not favour either system. The author had seen some of the
original grades before, so this is not a fully independent annotation.

For the eight claim-verification questions we additionally fixed a gold
verdict (SUPPORTED, REFUTED or INSUFFICIENT) from the ORKG records, with the
set of verdicts accepted as correct; one question whose evidence is mixed is
excluded from verdict scoring.

\subsection{Retrieval Metrics}

Each system returns 8 papers. We report nDCG@8 with linear gains (2/1/0) and the
ideal ranking computed from all judged papers of the question, together with
precision, recall of grade-2 papers and hit rate at 8. Differences between
systems are tested with a paired bootstrap over questions (5,000 resamples,
95\,\% percentile intervals).

Graph retrieval depends on the LLM that classifies the question and extracts
entities. We therefore report it in two modes: \emph{oracle}, with the gold
question type and hand-written entities, which isolates the retrieval
strategies, and \emph{end-to-end}, with entities extracted by \texttt{qwen2.5:7b}.
Dense retrieval does not use these steps.

\subsection{Answer Metrics without an LLM Judge}

Answers follow a fixed template and name papers by their exact title. A context
paper counts as \emph{cited} when its title appears in the answer body; the
final list of sources is ignored, because models tend to copy all context
papers into it. Per answer we compute:

\begin{itemize}
\item \textbf{CitePrec}: mean gold grade of the cited papers divided by 2;
\item \textbf{CiteUtil}: share of the grade-2 papers in the context that the answer cites;
\item \textbf{False abstention}: the answer cites nothing (or returns INSUFFICIENT)
  although the context contains a grade-2 paper.
\end{itemize}

The \textbf{gold Score} of an answer is $F_1$(CitePrec, CiteUtil) when the
context contains a grade-2 paper; otherwise it is 1 for an abstention and
CitePrec (at most 0.5) for an answer. Claim questions are scored by the
verdict, where INSUFFICIENT is also accepted when
none of the deciding papers was retrieved. This score is appropriate for
comparing LLMs on the same retrieval, but it rewards abstention when retrieval
fails, so a weaker retriever can obtain a higher Score. To compare retrieval
systems we therefore use a \textbf{strict end-to-end score} (E2E): no credit for
abstention, claim questions scored only by the gold verdict, and $F_1$(CitePrec,
CiteUtil) otherwise.

\subsection{Claim-Level Faithfulness}

Following claim-decomposition approaches \cite{min2023factscore,es2024ragas}, a
local judge (\texttt{qwen3:8b} via Ollama, not one of the evaluated models)
splits each answer body into at most 25 atomic claims and checks every claim
separately against the full context the model saw, labelling it supported (S),
unsupported (U) or contradicted (C). For S and C the judge must quote the
deciding context line verbatim; if the quote is not found in the context the
label becomes U. Faithfulness is the share of S claims per answer, averaged per
model.

The judge was validated on 100 random claims labelled blind by a stronger
model, with all 21 disagreements reviewed by the author. Accuracy was 0.79
($\kappa = 0.51$) for S/U/C and 0.84 ($\kappa = 0.60$) for supported versus not
supported. Errors were balanced (7 too strict, 7 too lenient), so per-model
means are unbiased, but single labels are only moderately reliable. The judge
rarely uses C and labels misattributions as U, so we report U and C together as
\emph{unsupported}.

\subsection{LLM-Judge Baseline}

For comparison with common practice \cite{liu2023geval,zheng2023judging}, every
answer was also scored by an LLM judge on groundedness, relevance and
completeness (0--5). The judge receives the question, the answer, the titles of
the retrieved papers and answer hints, but not the context itself. It is the
same \texttt{qwen3:8b} model as the faithfulness judge, so differences between
the two reflect the evaluation protocol rather than the judge model.

\subsection{Experimental Setup}

The two retrieval systems share the whole pipeline except the choice of the 8
papers (Section~\ref{sec:system}): the same full ORKG records, context rows,
prompt and LLM. Dense retrieval embeds one text card per paper (title, research
field and contribution values, at most 1,500 characters) for all 68,806 papers
in the ORKG RDF dump of 18 October 2026 with \texttt{mxbai-embed-large}
\cite{mxbai2024}, and ranks by cosine similarity; to test sensitivity to the
embedder we repeat retrieval with \texttt{nomic-embed-text} \cite{nussbaum2024nomic}.
A reciprocal rank fusion (RRF, $k = 60$) \cite{cormack2009rrf} of the two ranked lists
is computed at the retrieval level only.

Nine LLMs answer with graph retrieval: five local models via Ollama
(\texttt{qwen2.5:7b}, \texttt{llama3}, \texttt{llama3.2}, \texttt{mistral},
\texttt{phi3:mini}; context window 8,192 tokens) and four cloud models via the
Groq free tier (\texttt{gpt-oss-20b}, \texttt{gpt-oss-120b}, \texttt{qwen3.8-27b},
\texttt{allam-2-7b}). The graph--dense comparison of answers uses one local and
one cloud model (\texttt{qwen2.5:7b} and \texttt{gpt-oss-20b}); \texttt{qwen2.5:7b}
was run twice with each retriever to measure run-to-run variance.
